\documentclass[10pt]{article}
\usepackage[margin=0.85in]{geometry}
\usepackage{amsmath,amssymb,amsthm,booktabs,hyperref}
\newtheorem{lemma}{Lemma}
\newtheorem{theorem}{Theorem}
\newcommand{\SDR}{\mathrm{SDR}}
\title{Failure of contiguous-minor closure for Star of David matrices}
\author{Principia Math in Progress}
\date{13 September 2026}
\begin{document}
\maketitle

\paragraph{Task 1. Resolve Sun's contiguous-minor conjecture.}
For \(j\ge1\), Definition 2.12 of \cite{Sun} sets
\[
 A^{[j]}_{n,k}=\det[A_{n+i,k+\ell}]_{0\le i,\ell<j}.
\]
Conjecture 2.15 asserts that \(A\in\SDR_\infty\) implies
\(A^{[j]}\in\SDR_\infty\) for every \(j\ge1\).
We refute the assertion already for \(j=2\).
We use Sun's source definition: \(\SDR\) means \emph{Star of David rule}.
The added phrase ``sign-determined Riordan matrices'' in MathDB is
inconsistent with that definition. No Riordan-array restriction is
imposed in the source conjecture; our result resolves that original
assertion, and makes no assertion about a separately restricted class.

For an infinite lower triangular matrix \(X\), with zero entries above
the diagonal, put
\begin{align*}
 L_{p,r}(n,k;X)&=
 \prod_{i=0}^{r}X_{n+i,k+r-i}
 \prod_{i=0}^{p-r-1}X_{n+p-i,k+r+i+1},\\
 R_{p,r}(n,k;X)&=
 \prod_{i=0}^{r}X_{n+p-i,k+p-r+i}
 \prod_{i=0}^{p-r-1}X_{n+i,k+p-r-i-1}.
\end{align*}
Definition 1.1 in \cite{Sun} requires \(L_{p,r}=R_{p,r}\) for all
\(n\ge k\ge0\), \(2\le p\le m-1\), and \(0\le r\le p-1\).
Membership in \(\SDR_\infty\) means these identities hold for every
\(p\ge2\).

\begin{lemma}\label{extension}
Let \(M\) be a lower triangular \(d\times d\) matrix satisfying all
the displayed identities whose row indices lie in \(\{0,\ldots,d-1\}\).
Then \(M\oplus I\in\SDR_\infty\), where \(I\) is the infinite identity.
\end{lemma}
\begin{proof}
Each product uses every row \(n,\ldots,n+p\) and every column
\(k,\ldots,k+p\) exactly once. If \(n>k\) and \(n+p\ge d\),
row \(n+p\) belongs to the identity tail. Its only nonzero entry is
in column \(n+p\), exceeding all the columns in either product.
Consequently both products vanish. If \(n+p<d\), the assumed
finite identities apply.

It remains to handle \(n=k\). For \(r\ge1\), the first factor
\(X_{n,k+r}\) of \(L\) is above the diagonal; for \(r=0\),
the factor with \(i=p-1\) in its second product is
\(X_{n+1,k+p}=0\). In \(R\), if \(r\le p-2\), its second
product contains \(X_{n,k+p-r-1}=0\); if \(r=p-1\), its first
product contains \(X_{n+1,k+p}=0\). Thus both sides vanish here too.
\end{proof}

Consider
\[
 M=\begin{pmatrix}
 1&0&0&0&0\\1&1&0&0&0\\0&1&1&0&0\\0&0&1&1&0\\0&1&0&2&1
 \end{pmatrix},\qquad A=M\oplus I.
\]
All identities within \(M\) hold. To check this without any unbounded
computation, cases \(n=k\) were proved in the lemma. For \(n>k\)
and \(n+p\le4\), the only possibilities are the following:
\[
\begin{array}{cccll}
\toprule
p&(n,k)&r&L&R\\
\midrule
2&(1,0)&0&1\cdot0\cdot1=0&1\cdot1\cdot0=0\\
2&(2,0)&0&0\cdot1\cdot1=0&0\cdot1\cdot0=0\\
2&(2,1)&0&1\cdot0\cdot1=0&2\cdot1\cdot0=0\\
3&(1,0)&0&0\quad(M_{2,3}=0)&0\quad(M_{1,2}=0)\\
3&(1,0)&1&0\quad(M_{2,0}=0)&0\quad(M_{2,0}=0)\\
\bottomrule
\end{array}
\]
For \(p=2,r=1\), the two sides are the \(r=0\) sides interchanged.
For \(p=3,r=2\), they are the \(r=0\) sides interchanged.
There are no further cases with \(n>k\). The lemma proves
\(A\in\SDR_\infty\).

\begin{theorem}
For the infinite matrix \(A\) above, \(A\in\SDR_\infty\) but
\(A^{[2]}\notin\SDR_3\), and hence \(A^{[2]}\notin\SDR_\infty\).
\end{theorem}
\begin{proof}
Put \(C=A^{[2]}\). Its six entries needed for the order-three rule are
\[
 C_{1,0}=\det\begin{pmatrix}1&1\\0&1\end{pmatrix}=1,\qquad
 C_{3,1}=\det\begin{pmatrix}0&1\\1&0\end{pmatrix}=-1,\qquad
 C_{2,2}=1,
\]
\[
 C_{3,2}=\det\begin{pmatrix}1&1\\0&2\end{pmatrix}=2,\qquad
 C_{1,1}=1,\qquad
 C_{2,0}=\det\begin{pmatrix}0&1\\0&0\end{pmatrix}=0.
\]
Consequently the identity at \((p,n,k,r)=(2,1,0,0)\) fails:
\[
 L=C_{1,0}C_{3,1}C_{2,2}=-1
 \ne 0=C_{3,2}C_{1,1}C_{2,0}=R.
\]
Every entry used is a minor of the displayed five-by-five block;
the extension lemma verifies the stronger infinite-order hypothesis
on the input array, completing the counterexample.
\end{proof}

\paragraph{Verification and scope.}
The supplied script checks all 25 finite identities, the displayed
matrix operations, and the exact failed identity. It also checks a
larger truncation as a diagnostic. Lemma~\ref{extension} supplies the
proof for the entire infinite matrix; the diagnostic is not used to
infer an infinite assertion. Zero entries are allowed in Definition 1.1
and in the selected conjecture. The source explicitly imposes a
nonzero-entry hypothesis in other theorems, including Theorem 2.10,
but not here.
Remark 2.16 in the source gives an example with input only in order three. That example does not meet the infinite-order hypothesis here. Our input does. The same input also refutes the different inverse-closure conjecture, treated as MathDB 317353; the present determinant calculation resolves the separately stated minor operation.

\paragraph{Task 2. Minimize the minor counterexample.}
Find the smallest leading block supporting a failure of contiguous
two-by-two-minor closure under the infinite-order hypothesis.

\paragraph{Task 3. Find sufficient support conditions.}
Characterize zero patterns for which every contiguous-minor transform
preserves the Star of David identities.

\paragraph{Task 4. Study iterated minor transforms.}
Determine how failures propagate under repeated contiguous-minor
transforms, starting from sparse unit lower triangular arrays.

\paragraph{Task 5. Separate closure properties.}
Construct arrays that distinguish inverse closure from contiguous-minor
closure, and determine any implications between these two properties.

\begin{thebibliography}{1}
\bibitem{Sun} Y.~Sun, \emph{The Star of David rule}, Linear Algebra
Appl. \textbf{429} (2008), 1954--1961,
\href{https://doi.org/10.1016/j.laa.2008.05.027}{doi:10.1016/j.laa.2008.05.027};
\href{https://arxiv.org/abs/0805.1277}{arXiv:0805.1277v1}.
Source references in this note use the actual seven-page arXiv PDF.
\end{thebibliography}
\end{document}
